% ============================================================================
% CSCE 763: Digital Image Processing - Midterm Study Guide
% ============================================================================

\documentclass[10pt]{article}

% ============================================================================
% PACKAGE IMPORTS
% ============================================================================
\usepackage{amsmath}
\usepackage{amssymb}
\usepackage{tikz}
\usepackage[margin=0.75in]{geometry}
\usepackage{xcolor}
\usepackage{enumitem}
\usepackage{tcolorbox}
\usepackage{booktabs}
\usepackage{array}
\usepackage{multicol}
\usepackage{fancyhdr}

% ============================================================================
% PAGE SETUP
% ============================================================================
\pagestyle{fancy}
\setlength{\headheight}{14.5pt}
\fancyhf{}
\fancyhead[L]{CSCE 763 Midterm Study Guide}
\fancyhead[R]{Page \thepage}
\renewcommand{\headrulewidth}{0.4pt}

% ============================================================================
% COLORS
% ============================================================================
\definecolor{USC_Garnet}{HTML}{73000A}
\definecolor{USC_Atlantic}{HTML}{466A9F}
\definecolor{USC_Honeycomb}{HTML}{65780B}
\definecolor{USC_Congaree}{HTML}{1F414D}

% ============================================================================
% CUSTOM BOXES
% ============================================================================
\newtcolorbox{conceptbox}[1][]{
  colback=white,
  colframe=USC_Garnet,
  fonttitle=\bfseries,
  title=#1,
  sharp corners,
  colbacktitle=USC_Garnet,
  coltitle=white,
  left=3pt, right=3pt, top=2pt, bottom=2pt,
  boxsep=2pt,
}

\newtcolorbox{formulabox}{
  colback=USC_Atlantic!5,
  colframe=USC_Atlantic,
  sharp corners,
  left=3pt, right=3pt, top=2pt, bottom=2pt,
  boxsep=2pt,
}

\newtcolorbox{examplebox}{
  colback=USC_Honeycomb!5,
  colframe=USC_Honeycomb,
  sharp corners,
  left=3pt, right=3pt, top=2pt, bottom=2pt,
  boxsep=2pt,
}

% ============================================================================
% BEGIN DOCUMENT
% ============================================================================

\begin{document}

\begin{center}
{\LARGE\bfseries CSCE 763: Digital Image Processing}\\[0.3em]
{\Large Midterm Exam Study Guide}\\[0.2em]
{\small Organized by: Concept $\to$ Formula $\to$ Example $\to$ Where It Appears}
\end{center}

\vspace{0.3cm}

% ============================================================================
% SECTION 1: HUMAN VISION
% ============================================================================

\begin{conceptbox}[1. Human Vision \& Image Formation]

\textbf{Concept}: The eye forms images using a lens system. Similar triangles relate object size to image size.

\begin{formulabox}
\textbf{Lens Equation (Similar Triangles)}:
\[
\frac{h_{\text{image}}}{f} = \frac{h_{\text{object}}}{d} \quad \Rightarrow \quad h_{\text{image}} = \frac{f \cdot h_{\text{object}}}{d}
\]
where $f$ = focal length, $d$ = distance to object.

\textbf{Photoreceptors}:
\begin{itemize}[leftmargin=1.5em, itemsep=0pt, topsep=2pt]
    \item \textbf{Cones}: 6--7 million, color vision, concentrated in fovea, daylight (photopic)
    \item \textbf{Rods}: 75--150 million, grayscale, peripheral, low-light (scotopic)
\end{itemize}
\end{formulabox}

\begin{examplebox}
\textbf{Example}: Camera with $f = 35$ mm photographs 50 m tall building at 200 m distance.
\[
h_{\text{image}} = \frac{35 \text{ mm} \times 50{,}000 \text{ mm}}{200{,}000 \text{ mm}} = 8.75 \text{ mm}
\]
\end{examplebox}

\textbf{Appears in}: Lecture 1--2, HW1 P1 (cone density)

\end{conceptbox}

% ============================================================================
% SECTION 1B: PERCEPTION PHENOMENA
% ============================================================================

\begin{conceptbox}[1b. Brightness Perception \& Optical Illusions]

\begin{formulabox}
\textbf{Weber Ratio}: $\frac{\Delta I_c}{I} \approx 0.02$ \hfill (JND $\approx$ 2\% of background)

Human eye can discriminate $\sim$50 intensity levels, but this varies with background.

\textbf{Mach Bands}: Perceived bright/dark bands at intensity edges that don't exist in the actual intensity profile. Caused by lateral inhibition in retina.

\textbf{Simultaneous Contrast}: Same gray appears lighter on dark background, darker on light background. The perceived intensity depends on surroundings.
\end{formulabox}

\begin{examplebox}
\textbf{Example}: At background intensity $I = 100$, the just-noticeable difference is:
\[
\Delta I_c = 0.02 \times 100 = 2
\]
You can distinguish $I = 100$ from $I = 102$, but not from $I = 101$.
\end{examplebox}

\textbf{Appears in}: Lecture 1--2

\end{conceptbox}

% ============================================================================
% SECTION 1C: RESOLUTION
% ============================================================================

\begin{conceptbox}[1c. Spatial \& Intensity Resolution]

\begin{formulabox}
\textbf{Spatial Resolution}: Number of pixels per unit length (dpi, pixels/mm).
\begin{itemize}[leftmargin=1.5em, itemsep=0pt, topsep=2pt]
    \item Higher spatial resolution $\to$ finer detail, larger file
    \item $M \times N$ image with $k$ bits/pixel requires $M \times N \times k$ bits storage
\end{itemize}

\textbf{Intensity Resolution}: Number of bits per pixel ($k$ bits $\to 2^k$ gray levels).
\begin{itemize}[leftmargin=1.5em, itemsep=0pt, topsep=2pt]
    \item $k = 8$ bits $\to$ 256 levels (standard grayscale)
    \item Low bit depth causes \textbf{false contouring} (visible bands in smooth gradients)
\end{itemize}

\textbf{Isopreference Curves}: Tradeoff between spatial and intensity resolution for subjective quality.
\end{formulabox}

\begin{examplebox}
\textbf{Example}: A $512 \times 512$ image with 8 bits/pixel requires:
\[
512 \times 512 \times 8 = 2{,}097{,}152 \text{ bits} = 256 \text{ KB}
\]
Reducing to 4 bits/pixel halves storage but may show false contouring.
\end{examplebox}

\textbf{Appears in}: Lecture 2

\end{conceptbox}

% ============================================================================
% SECTION 1D: INTERPOLATION
% ============================================================================

\begin{conceptbox}[1d. Interpolation Methods]

\textbf{Concept}: Used for image resizing, rotation, and geometric transformations.

\begin{formulabox}
\textbf{Nearest Neighbor}: Use value of closest pixel. Fast but causes blocky artifacts.

\textbf{Bilinear Interpolation}: Weighted average of 4 nearest neighbors.
\[
v(x,y) = ax + by + cxy + d
\]
Smoother than nearest neighbor, slight blurring.

\textbf{Bicubic Interpolation}: Uses 16 neighbors (4$\times$4 grid).
\[
v(x,y) = \sum_{i=0}^{3}\sum_{j=0}^{3} a_{ij} x^i y^j
\]
Best quality, most computationally expensive.
\end{formulabox}

\begin{examplebox}
\textbf{Quality comparison} (for enlargement):
\begin{itemize}[leftmargin=1.5em, itemsep=0pt, topsep=2pt]
    \item Nearest neighbor: pixelated edges
    \item Bilinear: smooth but slightly blurry
    \item Bicubic: sharpest edges, best for photographs
\end{itemize}
\end{examplebox}

\textbf{Appears in}: Lecture 2--3

\end{conceptbox}

% ============================================================================
% SECTION 2: PIXEL RELATIONSHIPS
% ============================================================================

\begin{conceptbox}[2. Pixel Neighborhoods \& Adjacency]

\textbf{Concept}: Pixels can be adjacent in different ways depending on which neighbors count.

\begin{formulabox}
\textbf{4-Neighbors} $N_4(p)$: pixels at $(x\pm 1, y)$ and $(x, y\pm 1)$ --- orthogonal only

\textbf{8-Neighbors} $N_8(p)$: 4-neighbors plus diagonals $(x\pm 1, y\pm 1)$

\textbf{m-Adjacency}: Two pixels $p, q$ with values in $V$ are m-adjacent if:
\begin{enumerate}[leftmargin=1.5em, itemsep=0pt, topsep=2pt]
    \item $q \in N_4(p)$, OR
    \item $q \in N_D(p)$ AND $N_4(p) \cap N_4(q) \cap V = \emptyset$ (no shared 4-neighbor in $V$)
\end{enumerate}
\end{formulabox}

\begin{examplebox}
\textbf{Example}: For $V = \{1\}$, pixels at $(0,0)=1$ and $(1,1)=1$ are m-adjacent only if neither $(0,1)$ nor $(1,0)$ equals 1. This eliminates path ambiguity.
\end{examplebox}

\textbf{Appears in}: Lecture 2--3, HW1 P2 (adjacency), HW1 P3 (paths), Midterm P1a

\end{conceptbox}

% ============================================================================
% SECTION 3: DISTANCE METRICS
% ============================================================================

\begin{conceptbox}[3. Distance Metrics]

\begin{formulabox}
\textbf{Euclidean}: $D_e(p,q) = \sqrt{(x_1-x_2)^2 + (y_1-y_2)^2}$

\textbf{City-Block (Manhattan)}: $D_4(p,q) = |x_1-x_2| + |y_1-y_2|$

\textbf{Chessboard}: $D_8(p,q) = \max(|x_1-x_2|, |y_1-y_2|)$
\end{formulabox}

\begin{examplebox}
\textbf{Example}: From $p=(0,0)$ to $q=(4,4)$:
\begin{itemize}[leftmargin=1.5em, itemsep=0pt, topsep=2pt]
    \item $D_e = \sqrt{32} = 4\sqrt{2} \approx 5.66$
    \item $D_4 = 4 + 4 = 8$
    \item $D_8 = \max(4,4) = 4$
\end{itemize}
\end{examplebox}

\textbf{Appears in}: Lecture 2--3, HW1 P3 (path lengths)

\end{conceptbox}

% ============================================================================
% SECTION 4: SET OPERATIONS
% ============================================================================

\begin{conceptbox}[4. Set Operations on Images]

\begin{formulabox}
\textbf{Union}: $A \cup B$ --- pixels in $A$ OR $B$

\textbf{Intersection}: $A \cap B$ --- pixels in $A$ AND $B$

\textbf{Complement}: $A^c$ --- pixels NOT in $A$

\textbf{Difference}: $A - B = A \cap B^c$ --- pixels in $A$ but not $B$
\end{formulabox}

\begin{examplebox}
\textbf{Example}: Shaded region where exactly $A$ and $D$ overlap but not $B$ or $C$:
\[
(A \cap D \cap B^c \cap C^c)
\]
\end{examplebox}

\textbf{Appears in}: Lecture 2--3, HW1 P4 (set expressions), Midterm P2

\end{conceptbox}

% ============================================================================
% SECTION 5: LINEAR VS NONLINEAR OPERATORS
% ============================================================================

\begin{conceptbox}[5. Linear vs Nonlinear Operators]

\begin{formulabox}
\textbf{Linear Operator}: $H$ is linear if $H[a_1 f_1 + a_2 f_2] = a_1 H[f_1] + a_2 H[f_2]$

\textbf{Examples}:
\begin{itemize}[leftmargin=1.5em, itemsep=0pt, topsep=2pt]
    \item \textbf{Linear}: Averaging, Laplacian, convolution, correlation
    \item \textbf{Nonlinear}: Median, max, min, log transform, gamma correction
\end{itemize}
\end{formulabox}

\begin{examplebox}
\textbf{Example}: Prove median is nonlinear. Let $f_1 = [1,2,3]$, $f_2 = [3,2,1]$.

$\text{median}(f_1) = 2$, $\text{median}(f_2) = 2$, so $\text{median}(f_1) + \text{median}(f_2) = 4$

But $\text{median}(f_1 + f_2) = \text{median}([4,4,4]) = 4 \ne 4$... Actually works here.

Try $f_1 = [1,5,2]$, $f_2 = [2,1,5]$: median$(f_1)=2$, median$(f_2)=2$, sum $=4$

$f_1 + f_2 = [3,6,7]$, median $= 6 \ne 4$ \quad $\Rightarrow$ Nonlinear!
\end{examplebox}

\textbf{Appears in}: Lecture 6--7, HW2 P1 (median nonlinearity)

\end{conceptbox}

% ============================================================================
% SECTION 6: IMAGE ARITHMETIC
% ============================================================================

\begin{conceptbox}[6. Image Arithmetic Operations]

\begin{formulabox}
\textbf{Addition}: $g(x,y) = f_1(x,y) + f_2(x,y)$ \hfill (noise reduction by averaging)

\textbf{Subtraction}: $g(x,y) = f_1(x,y) - f_2(x,y)$ \hfill (change detection)

\textbf{Multiplication}: $g(x,y) = f_1(x,y) \times f_2(x,y)$ \hfill (masking, ROI selection)
\end{formulabox}

\begin{examplebox}
\textbf{Example}: Defect detection in manufacturing.

Store ``golden'' reference image $f_{\text{ref}}$. For each new product image $f_{\text{new}}$:
\[
\text{defect} = |f_{\text{new}} - f_{\text{ref}}|
\]
Nonzero regions indicate missing components. \textbf{Requirements}: consistent lighting, alignment, no vibration.
\end{examplebox}

\textbf{Appears in}: Lecture 3--4, HW2 P2 (defect detection)

\end{conceptbox}

% ============================================================================
% SECTION 6B: IMAGE AVERAGING FOR NOISE REDUCTION
% ============================================================================

\begin{conceptbox}[6b. Image Averaging for Noise Reduction]

\begin{formulabox}
\textbf{Model}: $g_i(x,y) = f(x,y) + \eta_i(x,y)$ where $\eta$ is noise with variance $\sigma_\eta^2$

\textbf{Average of $K$ noisy images}:
\[
\bar{g}(x,y) = \frac{1}{K}\sum_{i=1}^{K} g_i(x,y)
\]

\textbf{Key Result}: The variance of the averaged image is:
\[
\boxed{\sigma_{\bar{g}}^2 = \frac{1}{K}\sigma_\eta^2}
\]

\textbf{Standard deviation}: $\sigma_{\bar{g}} = \frac{\sigma_\eta}{\sqrt{K}}$

As $K \to \infty$, $\bar{g}(x,y) \to f(x,y)$ (noise averages to zero).
\end{formulabox}

\begin{examplebox}
\textbf{Example}: To reduce noise standard deviation by half, you need:
\[
\frac{\sigma_\eta}{\sqrt{K}} = \frac{\sigma_\eta}{2} \quad \Rightarrow \quad K = 4 \text{ images}
\]
To reduce by factor of 10, need $K = 100$ images.
\end{examplebox}

\textbf{Appears in}: Lecture 3--4

\end{conceptbox}

% ============================================================================
% SECTION 7: INTENSITY TRANSFORMATIONS
% ============================================================================

\begin{conceptbox}[7. Intensity Transformations]

\begin{formulabox}
\textbf{Log Transform}: $s = c \cdot \log(1 + r)$ \hfill (compress dynamic range, brighten darks)

\textbf{Gamma (Power-Law)}: $s = c \cdot r^\gamma$
\begin{itemize}[leftmargin=1.5em, itemsep=0pt, topsep=2pt]
    \item $\gamma < 1$: brightens dark regions (expand darks)
    \item $\gamma > 1$: darkens image (compress darks)
    \item $\gamma = 1$: identity
\end{itemize}

\textbf{Contrast Stretching}: $s = \frac{r - r_{\min}}{r_{\max} - r_{\min}} \cdot (L-1)$
\end{formulabox}

\begin{examplebox}
\textbf{Example}: For $\gamma = 0.4$ on a dark X-ray, dark pixels ($r=0.1$) map to:
\[
s = (0.1)^{0.4} \approx 0.40
\]
Brightens the image significantly.
\end{examplebox}

\textbf{Appears in}: Lecture 3--5, Midterm P1b (fireworks at night)

\end{conceptbox}

% ============================================================================
% SECTION 6: HISTOGRAM PROCESSING
% ============================================================================

\begin{conceptbox}[8. Histogram Equalization]

\begin{formulabox}
\textbf{Normalized Histogram}: $p_r(r_k) = \frac{n_k}{MN}$ where $n_k$ = pixel count at level $k$

\textbf{Histogram Equalization}:
\[
s_k = T(r_k) = (L-1) \sum_{i=0}^{k} p_r(r_i) = (L-1) \cdot \text{CDF}(r_k)
\]
Round $s_k$ to nearest integer.
\end{formulabox}

\begin{examplebox}
\textbf{Example}: 8-level image with $p_r = [0.1, 0.15, 0.25, 0.2, 0.125, 0.075, 0.0625, 0.0375]$

CDF: $[0.1, 0.25, 0.5, 0.7, 0.825, 0.9, 0.9625, 1.0]$

$s_k = 7 \times \text{CDF}$: $[0.7, 1.75, 3.5, 4.9, 5.775, 6.3, 6.7375, 7]$

Rounded: $[1, 2, 4, 5, 6, 6, 7, 7]$
\end{examplebox}

\textbf{Appears in}: Lecture 5--6, HW2 P3 (idempotency), HW2 P4 (discrete matching)

\end{conceptbox}

% ============================================================================
% SECTION 7: HISTOGRAM SPECIFICATION
% ============================================================================

\begin{conceptbox}[9. Histogram Specification (Matching)]

\begin{formulabox}
\textbf{Goal}: Transform image with PDF $p_r(r)$ to have PDF $p_z(z)$

\textbf{Method} (use equalization as a bridge):
\begin{enumerate}[leftmargin=1.5em, itemsep=0pt, topsep=2pt]
    \item Equalize input: $s = T_r(r) = (L-1)\int_0^r p_r(w)\,dw$
    \item Equalize desired: $s = T_z(z) = (L-1)\int_0^z p_z(w)\,dw$
    \item Set equal and solve: $T_r(r) = T_z(z) \Rightarrow z = T_z^{-1}(T_r(r))$
\end{enumerate}
\end{formulabox}

\begin{examplebox}
\textbf{Example}: Given $p_r(r) = \frac{4r}{3(L-1)^2}$ and $p_z(z) = \frac{2z}{(L-1)^2}$

$T_r(r) = \frac{2r^2}{3(L-1)}$, \quad $T_z(z) = \frac{z^2}{L-1}$

Setting equal: $\frac{2r^2}{3(L-1)} = \frac{z^2}{L-1} \Rightarrow z = \sqrt{\frac{2}{3}}\,r$
\end{examplebox}

\textbf{Appears in}: Lecture 5--6, HW2 P5 (continuous), Midterm P4

\end{conceptbox}

% ============================================================================
% SECTION 8: SPATIAL FILTERING
% ============================================================================

\begin{conceptbox}[10. Spatial Filtering --- Smoothing]

\begin{formulabox}
\textbf{Averaging Filter}: $\frac{1}{9}\begin{bmatrix} 1 & 1 & 1 \\ 1 & 1 & 1 \\ 1 & 1 & 1 \end{bmatrix}$ \hfill (blur, reduce noise)

\textbf{Gaussian Filter}: $G(x,y) = \frac{1}{2\pi\sigma^2}e^{-(x^2+y^2)/2\sigma^2}$ \hfill (smooth blur)

\textbf{Median Filter}: Output = median of neighborhood \hfill (removes salt-and-pepper)
\end{formulabox}

\begin{examplebox}
\textbf{Example}: Neighborhood $[15, 20, 255, 18, 22, 19, 21, 17, 0]$

Arithmetic mean: $(15+20+255+18+22+19+21+17+0)/9 = 43$

Median: Sort $\to [0,15,17,18,19,20,21,22,255]$, median $= 19$
\end{examplebox}

\textbf{Appears in}: Lecture 6--7, HW3 P1 (blurring histograms), HW3 P3 (bar merging)

\end{conceptbox}

% ============================================================================
% SECTION 9: SHARPENING FILTERS
% ============================================================================

\begin{conceptbox}[11. Spatial Filtering --- Sharpening]

\begin{formulabox}
\textbf{Laplacian} (4-neighbor): $\begin{bmatrix} 0 & -1 & 0 \\ -1 & 4 & -1 \\ 0 & -1 & 0 \end{bmatrix}$ \quad (8-neighbor): $\begin{bmatrix} -1 & -1 & -1 \\ -1 & 8 & -1 \\ -1 & -1 & -1 \end{bmatrix}$

\textbf{Sharpened Image}: $g(x,y) = f(x,y) - \nabla^2 f(x,y)$ \hfill (subtract Laplacian)

\textbf{High-Boost}: $g = A \cdot f - \text{lowpass}(f) = (A-1)f + \text{highpass}(f)$

\textbf{Sobel} (horizontal edges): $\begin{bmatrix} -1 & -2 & -1 \\ 0 & 0 & 0 \\ 1 & 2 & 1 \end{bmatrix}$ \quad (vertical): $\begin{bmatrix} -1 & 0 & 1 \\ -2 & 0 & 2 \\ -1 & 0 & 1 \end{bmatrix}$
\end{formulabox}

\begin{examplebox}
\textbf{Example}: Filter $\frac{1}{9}\begin{bmatrix} -1 & -1 & -1 \\ -1 & -17 & -1 \\ -1 & -1 & -1 \end{bmatrix}$

This is a sharpening filter with inversion (all negative). Center weight $\gg$ neighbors $\Rightarrow$ edge enhancement.
\end{examplebox}

\textbf{Appears in}: Lecture 6--7, HW4 P1 (highboost k=3), Midterm P3

\end{conceptbox}

% ============================================================================
% SECTION 10: CONVOLUTION
% ============================================================================

\begin{conceptbox}[12. Convolution vs Correlation]

\begin{formulabox}
\textbf{Correlation}: $w(x,y) \star f(x,y) = \sum_s \sum_t w(s,t) \cdot f(x+s, y+t)$

\textbf{Convolution}: $w(x,y) * f(x,y) = \sum_s \sum_t w(s,t) \cdot f(x-s, y-t)$

\textbf{Key difference}: Convolution flips the filter 180°. For symmetric filters, they're the same.

\textbf{Property}: If filter coefficients sum to zero, output sum is also zero.

\textbf{Zero Padding}: Extend image with zeros at boundaries for full output.
\end{formulabox}

\begin{examplebox}
\textbf{Example}: Position $(0,0)$ with zero-padded neighborhood $\begin{bmatrix} 0 & 0 & 0 \\ 0 & 3 & 5 \\ 0 & 7 & 18 \end{bmatrix}$

Filter $\frac{1}{9}\begin{bmatrix} -1 & -1 & -1 \\ -1 & -17 & -1 \\ -1 & -1 & -1 \end{bmatrix}$:

$\frac{1}{9}[(-1)(0+0+0+0+5+0+7+18) + (-17)(3)] = \frac{1}{9}[-30-51] = -9$
\end{examplebox}

\textbf{Appears in}: Lecture 6--7, HW3 P2 (properties), HW3 P5 (full output), Midterm P3b

\end{conceptbox}

% ============================================================================
% SECTION 11: FOURIER TRANSFORM
% ============================================================================

\begin{conceptbox}[13. Fourier Transform]

\begin{formulabox}
\textbf{Continuous FT}: $F(u) = \int_{-\infty}^{\infty} f(t)\, e^{-j2\pi ut}\,dt$

\textbf{Rectangular Pulse Formula} (height $h$, from $a$ to $b$):
\[
\boxed{\int_a^b h\, e^{-j2\pi ut}\,dt = \frac{h}{j2\pi u}\left(e^{-j2\pi ua} - e^{-j2\pi ub}\right)}
\]

\textbf{Sifting Property}: $\int_{-\infty}^{\infty} f(t)\,\delta(t-t_0)\,dt = f(t_0)$
\end{formulabox}

\begin{examplebox}
\textbf{Example}: $\int_{-\infty}^{\infty} (t^2 + 3t + 5)\,\delta(t-2)\,dt = (2)^2 + 3(2) + 5 = 15$
\end{examplebox}

\textbf{Appears in}: Lecture 7--8, HW4 P2 (piecewise FT), Midterm Bonus

\end{conceptbox}

% ============================================================================
% SECTION 12: DFT
% ============================================================================

\begin{conceptbox}[14. Discrete Fourier Transform]

\begin{formulabox}
\textbf{1D DFT}: $F(u) = \sum_{x=0}^{M-1} f(x)\, e^{-j2\pi ux/M}$

\textbf{Inverse}: $f(x) = \frac{1}{M}\sum_{u=0}^{M-1} F(u)\, e^{j2\pi ux/M}$

\textbf{2D DFT}: $F(u,v) = \sum_{x=0}^{M-1}\sum_{y=0}^{N-1} f(x,y)\, e^{-j2\pi(ux/M + vy/N)}$

\textbf{Convolution Theorem}: $f * h \leftrightarrow F \cdot H$ \hfill (convolution $\to$ multiplication)
\end{formulabox}

\textbf{Appears in}: Lecture 7--8

\end{conceptbox}

% ============================================================================
% SECTION 13: FREQUENCY DOMAIN FILTERS
% ============================================================================

\begin{conceptbox}[15. Frequency Domain Filters]

\begin{formulabox}
\textbf{Ideal Lowpass}: $H(u,v) = \begin{cases} 1 & D(u,v) \le D_0 \\ 0 & D(u,v) > D_0 \end{cases}$ \hfill (causes ringing)

\textbf{Butterworth Lowpass}: $H(u,v) = \frac{1}{1 + [D(u,v)/D_0]^{2n}}$ \hfill (smooth rolloff)

\textbf{Gaussian Lowpass}: $H(u,v) = e^{-D^2(u,v)/(2D_0^2)}$ \hfill (no ringing)

\textbf{Highpass}: $H_{\text{HP}} = 1 - H_{\text{LP}}$

where $D(u,v) = \sqrt{(u-M/2)^2 + (v-N/2)^2}$
\end{formulabox}

\textbf{Appears in}: Lecture 8+

\end{conceptbox}

% ============================================================================
% SECTION 14: NOISE MODELS
% ============================================================================

\begin{conceptbox}[16. Noise Models]

\begin{formulabox}
\textbf{Gaussian}: $p(z) = \frac{1}{\sqrt{2\pi}\sigma}\exp\left[-\frac{(z-\mu)^2}{2\sigma^2}\right]$

\textbf{Salt-and-Pepper}: $p(z) = \begin{cases} P_a & z = a \text{ (black)} \\ P_b & z = b \text{ (white)} \\ 0 & \text{otherwise} \end{cases}$

\textbf{Uniform}: $p(z) = \frac{1}{b-a}$ for $a \le z \le b$
\end{formulabox}

\begin{examplebox}
\textbf{Best filters for each noise type}:
\begin{itemize}[leftmargin=1.5em, itemsep=0pt, topsep=2pt]
    \item Gaussian noise $\to$ Arithmetic mean, Gaussian smoothing
    \item Salt-and-pepper $\to$ Median filter
    \item Uniform noise $\to$ Arithmetic mean
\end{itemize}
\end{examplebox}

\textbf{Appears in}: Lecture 8+, HW4 P3 (10 filter types)

\end{conceptbox}

% ============================================================================
% SECTION 15: RESTORATION FILTERS
% ============================================================================

\begin{conceptbox}[17. Restoration Filters (Mean Filters)]

\begin{formulabox}
\textbf{Arithmetic Mean}: $\hat{f} = \frac{1}{mn}\sum g(s,t)$

\textbf{Geometric Mean}: $\hat{f} = \left[\prod g(s,t)\right]^{1/mn}$ \hfill (less blurring than arithmetic)

\textbf{Harmonic Mean}: $\hat{f} = \frac{mn}{\sum 1/g(s,t)}$ \hfill (good for salt noise)

\textbf{Contraharmonic Mean}: $\hat{f} = \frac{\sum g(s,t)^{Q+1}}{\sum g(s,t)^Q}$ \hfill ($Q>0$: pepper, $Q<0$: salt, $Q=0$: arithmetic)

\textbf{Median}: $\hat{f} = \text{median}\{g(s,t)\}$ \hfill (best for salt-and-pepper)

\textbf{Max}: $\hat{f} = \max\{g(s,t)\}$ \hfill (removes pepper/dark noise)

\textbf{Min}: $\hat{f} = \min\{g(s,t)\}$ \hfill (removes salt/bright noise)

\textbf{Midpoint}: $\hat{f} = \frac{1}{2}(\max + \min)$ \hfill (Gaussian + uniform noise)

\textbf{Alpha-Trimmed Mean}: Delete $d/2$ lowest and $d/2$ highest, average rest
\end{formulabox}

\begin{examplebox}
\textbf{Example}: Neighborhood $[15, 20, 255, 18, 0, 22, 19, 21, 17]$, $d=2$

Sorted: $[0, 15, 17, 18, 19, 20, 21, 22, 255]$

Trim 1 lowest (0) and 1 highest (255): $[15, 17, 18, 19, 20, 21, 22]$

Alpha-trimmed mean: $(15+17+18+19+20+21+22)/7 = 18.86$
\end{examplebox}

\textbf{Appears in}: Lecture 8+, HW4 P3 (all 10 filter types)

\end{conceptbox}

% ============================================================================
% SECTION 18: MOTION BLUR AND WIENER FILTER
% ============================================================================

\begin{conceptbox}[18. Motion Blur \& Wiener Filter]

\begin{formulabox}
\textbf{Motion Blur Model}: $H(u,v) = \int_0^T e^{-j2\pi[u\,x_0(t) + v\,y_0(t)]}\,dt$

For uniform horizontal motion $x_0(t) = at/T$:
\[
H(u,v) = \frac{T}{\pi u a}\sin(\pi u a)\,e^{-j\pi u a}
\]

\textbf{Wiener Filter}: $\hat{F}(u,v) = \frac{H^*(u,v)}{|H(u,v)|^2 + K} \cdot G(u,v)$

where $K = S_\eta/S_f$ (noise-to-signal power ratio).
\end{formulabox}

\begin{examplebox}
\textbf{Example}: Given $H(u,v) = 2\pi\sigma^2 e^{-2\pi^2\sigma^2(u^2+v^2)}$ and ratio $K$:
\[
\hat{F} = \frac{2\pi\sigma^2 e^{-2\pi^2\sigma^2(u^2+v^2)}}{4\pi^2\sigma^4 e^{-4\pi^2\sigma^2(u^2+v^2)} + K} \cdot G
\]
\end{examplebox}

\textbf{Appears in}: Lecture 8+, HW4 P4 (motion blur), HW4 P5 (Wiener)

\end{conceptbox}

% ============================================================================
% SECTION 19: ORDER OF OPERATIONS
% ============================================================================

\begin{conceptbox}[19. Order of Linear Operations]

\begin{formulabox}
\textbf{Key Property}: Linear operations (convolutions) are \textbf{commutative}.

If $h_1$ and $h_2$ are linear filters:
\[
(f * h_1) * h_2 = (f * h_2) * h_1 = f * (h_1 * h_2)
\]

\textbf{Result}: Averaging then Laplacian = Laplacian then Averaging (same result)
\end{formulabox}

\begin{examplebox}
\textbf{Example}: Apply averaging (noise reduction) then Laplacian (edge detection).

Since both are linear, order doesn't matter mathematically. However, averaging first is often preferred in practice to reduce noise before edge detection.
\end{examplebox}

\textbf{Appears in}: Lecture 6--7, HW3 P4 (averaging + Laplacian order)

\end{conceptbox}

% ============================================================================
% QUICK REFERENCE: COMMON FILTERS
% ============================================================================

\vspace{0.3cm}
\begin{center}
\begin{tcolorbox}[colback=white, colframe=USC_Congaree, title=\textbf{Quick Reference: Common 3$\times$3 Filters}, sharp corners, colbacktitle=USC_Congaree, coltitle=white]

\begin{tabular}{@{}p{3.5cm}p{4cm}p{6cm}@{}}
\toprule
\textbf{Filter} & \textbf{Kernel} & \textbf{Effect} \\
\midrule
Averaging (box) & $\frac{1}{9}\begin{bmatrix} 1 & 1 & 1 \\ 1 & 1 & 1 \\ 1 & 1 & 1 \end{bmatrix}$ & Blur, noise reduction \\[1.5em]
Weighted avg. & $\frac{1}{16}\begin{bmatrix} 1 & 2 & 1 \\ 2 & 4 & 2 \\ 1 & 2 & 1 \end{bmatrix}$ & Smoother blur (Gaussian approx.) \\[1.5em]
Laplacian (4-conn) & $\begin{bmatrix} 0 & -1 & 0 \\ -1 & 4 & -1 \\ 0 & -1 & 0 \end{bmatrix}$ & Edge detection \\[1.5em]
Laplacian (8-conn) & $\begin{bmatrix} -1 & -1 & -1 \\ -1 & 8 & -1 \\ -1 & -1 & -1 \end{bmatrix}$ & Edge detection (diagonal) \\[1.5em]
Sobel $G_x$ & $\begin{bmatrix} -1 & 0 & 1 \\ -2 & 0 & 2 \\ -1 & 0 & 1 \end{bmatrix}$ & Vertical edge detection \\[1.5em]
Sobel $G_y$ & $\begin{bmatrix} -1 & -2 & -1 \\ 0 & 0 & 0 \\ 1 & 2 & 1 \end{bmatrix}$ & Horizontal edge detection \\[1.5em]
High-boost ($A=2$) & $\frac{1}{9}\begin{bmatrix} -1 & -1 & -1 \\ -1 & 17 & -1 \\ -1 & -1 & -1 \end{bmatrix}$ & Sharpen (enhance edges) \\[1.5em]
Prewitt $G_x$ & $\begin{bmatrix} -1 & 0 & 1 \\ -1 & 0 & 1 \\ -1 & 0 & 1 \end{bmatrix}$ & Vertical edge (simpler) \\
\bottomrule
\end{tabular}

\end{tcolorbox}
\end{center}

\end{document}
